\origpage{255}{GRAPH MINORS. VIII}
{\centering
{\Large Graph Minors. VIII. A Kuratowski Theorem for General Surfaces\par}\medskip

NEIL ROBERTSON\par
\textit{Department of Mathematics, Ohio State University, 231 West 18th Ave., Columbus, Ohio 43210}\par
AND\par
P. D. SEYMOUR$^*$\par
\textit{Bellcore, 445 South Street, Morristown, New Jersey 07960}\par
Received August 19, 1984; revised September 6, 1989\par
}
\begin{quote}
We prove that for any infinite set of graphs of bounded genus, some member of the set is isomorphic to a minor of another. As a consequence, for any surface $\Sigma$ there is a finite list of graphs, such that a general graph may be drawn in $\Sigma$ if an only if it topologically contains none of the graphs in the list.

\end{quote}

\sectionhead{1. \textsc{Introduction}}

Kuratowski's famous theorem of 1930 [6] asserts that a graph is planar if and only if it does not topologically contain $K_5$ or $K_{3,3}$. (Graphs are finite, and may have loops or multiple edges. A graph $G$ topologically contains $H$ if we may obtain a graph isomorphic to $H$ from some subgraph of $G$ by suppressing some divalent vertices.) There arose in the 1930's the proposal to find parallels to this theorem which apply to other surfaces.

If $\Sigma$ is a surface, let $T(\Sigma)$ denote the class of all graphs which cannot be drawn in $\Sigma$ and which are minimal with this property under topological containment. If $\Sigma$ is the plane or sphere then the members of $T(\Sigma)$ are precisely the graphs isomorphic to $K_5$ or $K_{3,3}$---this is another way to state Kuratowski's theorem. In general, a graph can be drawn in $\Sigma$ if and only if it topologically contains no members of $T(\Sigma)$. It remains then to determine $T(\Sigma)$ explicitly.

This appears to be very difficult, and there was little progress on the%
\origfootnote{Partially supported by NSF grant MCS 8103440.}
